% GENERATED FILE. Edit canonical lessons, then run npm run generate:books.
% canonical-source-hash: fnv1a64:e5b91f41148465a9
% canonical-lessons: RU-C142-nosit, RU-C142-derzhat, RU-C142-brosat, RU-C142-lovit, RU-C142-tolkat

\chapter{To carry, To hold, To throw, To catch, To push}
\label{ch:ru-to-carry-to-hold-to-throw-to-catch-to-push}
\hlchaptermodality{142}

\begin{chapteropening}
\textbf{By the end of this chapter:} \emph{I can say to carry, to hold, to throw, to catch and to push.}

Complete the last lesson of chapter 142: I can say to carry, to hold, to throw, to catch and to push.
\end{chapteropening}

\section[\texorpdfstring{\ru{носить} (nosít)}{носить (nosít)}]{\ru{носить} (nosít) — to carry; to wear}

\label{lesson:RU-C142-nosit}



\begin{quote}
Before the new one: say the Russian for to thank, then the Russian for to apologise.
\end{quote}

\subsection*{You'll want to know: \ru{носить}}
\textbf{\ru{носить}} (\emph{nosít}) — \textquotedblleft{}to carry; to wear\textquotedblright{}.

\begin{grammarlens}[title={What you've built}]
One more word for this chapter.
\end{grammarlens}

\subsection*{Your turn}
\begin{itemize}
\raggedright
  \item \emph{Say it:} \emph{nosít}
  \item \emph{Say it:} \emph{nosít}, once more
  \item \emph{Say it:} say \emph{nosít}
  \item \emph{Recall:} say the Russian for to thank, then the Russian for to apologise, then say \emph{nosít} again
\end{itemize}

\begin{tcolorbox}[breakable,colback=teal!4,colframe=teal!35!black,title={Before you move on}]
What does \ru{носить} mean? (\textquotedblleft{}To carry; to wear\textquotedblright{}.) Say it once more.
\end{tcolorbox}
\section[\texorpdfstring{\ru{держать} (derzhát)}{держать (derzhát)}]{\ru{держать} (derzhát) — to hold}

\label{lesson:RU-C142-derzhat}



\begin{quote}
Before the new one: say the Russian for to apologise, then the Russian for to carry; to wear.
\end{quote}

\subsection*{You'll want to know: \ru{держать}}
\textbf{\ru{держать}} (\emph{derzhát}) — \textquotedblleft{}to hold\textquotedblright{}.

\begin{grammarlens}[title={What you've built}]
One more word for this chapter.
\end{grammarlens}

\subsection*{Your turn}
\begin{itemize}
\raggedright
  \item \emph{Say it:} \emph{derzhát}
  \item \emph{Say it:} \emph{derzhát}, once more
  \item \emph{Say it:} say \emph{derzhát}
  \item \emph{Recall:} say the Russian for to apologise, then the Russian for to carry; to wear, then say \emph{derzhát} again
\end{itemize}

\begin{tcolorbox}[breakable,colback=teal!4,colframe=teal!35!black,title={Before you move on}]
What does \ru{держать} mean? (\textquotedblleft{}To hold\textquotedblright{}.) Say it once more.
\end{tcolorbox}
\section[\texorpdfstring{\ru{бросать} (brosát)}{бросать (brosát)}]{\ru{бросать} (brosát) — to throw}

\label{lesson:RU-C142-brosat}



\begin{quote}
Before the new one: say the Russian for to carry; to wear, then the Russian for to hold.
\end{quote}

\subsection*{You'll want to know: \ru{бросать}}
\textbf{\ru{бросать}} (\emph{brosát}) — \textquotedblleft{}to throw\textquotedblright{}.

\begin{grammarlens}[title={What you've built}]
One more word for this chapter.
\end{grammarlens}

\subsection*{Your turn}
\begin{itemize}
\raggedright
  \item \emph{Say it:} \emph{brosát}
  \item \emph{Say it:} \emph{brosát}, once more
  \item \emph{Say it:} say \emph{brosát}
  \item \emph{Recall:} say the Russian for to carry; to wear, then the Russian for to hold, then say \emph{brosát} again
\end{itemize}

\begin{tcolorbox}[breakable,colback=teal!4,colframe=teal!35!black,title={Before you move on}]
What does \ru{бросать} mean? (\textquotedblleft{}To throw\textquotedblright{}.) Say it once more.
\end{tcolorbox}
\section[\texorpdfstring{\ru{ловить} (lovít)}{ловить (lovít)}]{\ru{ловить} (lovít) — to catch}

\label{lesson:RU-C142-lovit}



\begin{quote}
Before the new one: say the Russian for to hold, then the Russian for to throw.
\end{quote}

\subsection*{You'll want to know: \ru{ловить}}
\textbf{\ru{ловить}} (\emph{lovít}) — \textquotedblleft{}to catch\textquotedblright{}.

\begin{grammarlens}[title={What you've built}]
One more word for this chapter.
\end{grammarlens}

\subsection*{Your turn}
\begin{itemize}
\raggedright
  \item \emph{Say it:} \emph{lovít}
  \item \emph{Say it:} \emph{lovít}, once more
  \item \emph{Say it:} say \emph{lovít}
  \item \emph{Recall:} say the Russian for to hold, then the Russian for to throw, then say \emph{lovít} again
\end{itemize}

\begin{tcolorbox}[breakable,colback=teal!4,colframe=teal!35!black,title={Before you move on}]
What does \ru{ловить} mean? (\textquotedblleft{}To catch\textquotedblright{}.) Say it once more.
\end{tcolorbox}
\section[\texorpdfstring{\ru{толкать} (tolkát)}{толкать (tolkát)}]{\ru{толкать} (tolkát) — to push}

\label{lesson:RU-C142-tolkat}



\begin{quote}
Before the new one: say the Russian for to throw, then the Russian for to catch.
\end{quote}

\subsection*{You'll want to know: \ru{толкать}}
\textbf{\ru{толкать}} (\emph{tolkát}) — \textquotedblleft{}to push\textquotedblright{}.

\begin{grammarlens}[title={What you've built}]
One more word for this chapter.
\end{grammarlens}

\subsection*{Your turn}
\begin{itemize}
\raggedright
  \item \emph{Say it:} \emph{tolkát}
  \item \emph{Say it:} \emph{tolkát}, once more
  \item \emph{Say it:} say \emph{tolkát}
  \item \emph{Recall:} say the Russian for to throw, then the Russian for to catch, then say \emph{tolkát} again
\end{itemize}

\begin{tcolorbox}[breakable,colback=teal!4,colframe=teal!35!black,title={Before you move on}]
What does \ru{толкать} mean? (\textquotedblleft{}To push\textquotedblright{}.) Say it once more.
\end{tcolorbox}
